\chapter{Как машина учится}
% полная версия: chapters/06_dofa.tex

Во всех схемах до сих пор силу связей подбирал инженер. В живом мозге инженера нет. Связи должны настраиваться сами, по ходу работы, --- и так, чтобы машина делала что-то полезное. Это и есть обучение.

Главное ограничение звучит очень строго. Каждое соединение меняет себя само, и знать оно может только то, что происходит рядом: работает ли отправитель, работает ли получатель и какие вещества до него доходят. Никто не может прислать соединению личную поправку. Это исключает главный метод обучения искусственных нейросетей, где ошибка пробегает сеть в обратную сторону и каждая связь получает свою собственную поправку.

\section*{Вместе --- значит связаны}

Самое простое правило: если отправитель и получатель работают вместе, связь между ними крепнет. Оно локально и не требует часов. Такое правило, если добавить к нему немного забывания, находит в потоке сигналов самое главное --- то, что меняется сильнее всего. Есть и версия, которая смотрит на время щелчков: связь крепнет, если отправитель щёлкнул \emph{перед} получателем, то есть мог быть причиной, и слабеет, если после.

Но у этих правил один предел: они учат тому, что \emph{часто}, а не тому, что \emph{полезно}. Соединение, которое видит только своих соседей, не знает, принесло ли действие сок или боль.

\section*{Третий сигнал}

Его приносит дофамин --- радиостанция «лучше, чем ожидалось». Пусть связь меняется, только когда совпали три вещи: работал отправитель, работал получатель и пришёл дофамин.

\pencilfig{6}{\pencilart{6}}{Связь крепнет, когда оба нейрона работали вместе и пришёл третий сигнал: «вышло лучше, чем ждали».}

Но вот трудность: награда приходит не сразу, а через секунду-другую. К этому моменту отправитель и получатель давно замолчали. Соединению нужна память о том, что оно участвовало. И она есть: совместная работа оставляет в соединении метку, как стикер, который сам отклеивается за секунду-две. Пришёл дофамин, пока стикер держится, --- связь меняется. Не пришёл --- метка исчезает бесследно. Это решает задачу адресации: дофамин один на всех, но меняются только те соединения, на которых висит стикер.

\section*{Шум как поиск}

Почему это правило ведёт к лучшему, а не куда попало? Представьте, что вы с завязанными глазами ищете вершину холма. Вы делаете случайный шаг, и кто-то кричит «выше!» или «ниже!». Если «выше» --- закрепляете шаг. Шаг за шагом вы поднимаетесь, ни разу не посмотрев на карту. Мозг делает то же самое: нейроны немного дрожат случайным образом, дофамин сообщает, стало ли лучше, и участвовавшие связи сдвигаются в выгодную сторону. Шум, который мешал передавать сообщения, здесь становится поиском.

Теперь понятно, почему дофамин сообщает не о награде, а о \emph{сюрпризе}. Если бы он просто кричал «сок!» после каждой удачи, он заодно закреплял бы всё, что сеть делала, --- и правильное, и неправильное. В нашей модели такая сеть действительно разгоняет свои связи в тысячи раз, а иногда навсегда застревает на худшем выборе. Вычитание ожидания делает обучение честным.

\section*{Цена одного провода}

У широковещательного обучения есть цена. Один сигнал на всех, а в нём смешан шум всех нейронов, влияющих на результат. Чем больше сеть, тем дольше каждому соединению выделять свою долю. Поэтому дофамин, по-видимому, учит только сравнительно небольшие цепи, выбирающие действия, а огромная кора учится в основном по местным правилам.

\section*{Кто считает сюрприз}

Остаётся вопрос: как дофаминовые нейроны вычисляют «лучше, чем ожидалось»? Нужен критик --- группа нейронов, которая всё время оценивает, сколько хорошего ещё впереди. Сюрприз --- это полученная награда плюс то, насколько подскочило ожидание. Зажглась лампочка --- ожидание подпрыгнуло, всплеск. Пришёл обещанный сок --- награда есть, но ожидание в тот же миг исчерпано, сюрприза нет. Сок не пришёл --- ожидание рухнуло впустую, провал. Все три случая из первой главы получаются сами собой. Что дофамин ведёт себя именно так, измерено. Как именно мозг вычисляет этот сюрприз --- пока гипотеза.

\fullversion{6}
